\section{緒論}

在當代資本主義的生產過程中，知識的性質與其在勞動中的角色日益受到學界重視。
\textcite{polanyiTacitDimension1966}提出「默會知識」（tacit knowledge）的概念，
指出人類的知識存在一種難以言說、無法完全編碼化的面向——「我們所知道的，
多於我們所能言說的」\footnote{%
  原文為 "We can know more than we can tell."，出自
  \citetitle{polanyiTacitDimension1966}，第4頁。}。
這一洞見對於理解當代勞動過程具有深遠意義。

近年來，生成式人工智慧（generative AI）的快速發展，對勞動市場造成顯著衝擊。
\textcite{brynjolfssonGenerativeAIWork2025}的研究顯示，AI 工具能夠顯著提升
客服工作者的生產力，平均提高約 15\%，且對低技能工作者的效益尤為明顯。
然而，此一現象是否意味著默會知識的重要性正在式微，仍有待進一步探討。
